% ------------------------------------------------------------------------
%
\begin{frame}{Naïve search}

Factor matching is common in text editing, although it is usually
better known as \textbf{exact string matching} in the academic field
of \emph{text algorithmics} (sometimes called \emph{stringology}).

\bigskip

Because of the asymmetric nature of the search, the word~\(p\) to find
will from now on be called the \textbf{pattern}, and the word~\(t\)
the \textbf{text}.

\bigskip

Assuming \(p\)~and~\(t\) are not empty, the naïve search simply
consists in checking whether \(p\)~is a prefix of~\(t\): if so, we are
done, otherwise, the pattern is shifted by one letter to the right
and, again, it is checked whether it is a prefix of that text suffix.

\bigskip

If \(p\)~cannot be shifted anymore because its end would surpass the
end of~\(t\), then it is not a factor: the search failed.

\end{frame}

% ------------------------------------------------------------------------
%
\begin{frame}{Naïve search}

\begin{figure}
\centering
\includegraphics[bb=75 621 352 715]{naive}
\end{figure}
Naïvely matching pattern~\(p\) against text~\(t\)
  (failure in grey), where \(\ind{p}{i} \neq \ind{t}{j}\) (the letters
\word{a}~and~\word{b} are not relevant in themselves).

\end{frame}

% ------------------------------------------------------------------------
%
\begin{frame}{Naïve search}

  An imperative program:

\begin{codebox}
\(\proc{Naive}(p, t)\)\\
\li \(i, j \gets 0, 0\)
\li \While \(i < \len{p}\) \LogAnd \(j < \len{t}\) \kw{do}
\li \> \If \(p[i] = t[j]\)
\li    \Then \(i, j \gets i + 1, j + 1\)
\li    \Else \(i, j \gets 0, j - i + 1\) \>\>\>\>\>\>\Comment We
slide \(p\) by 1.
       \End
\li \(\proc{Naive} \gets i = \len{p}\)
\end{codebox}

What test cases would be meaningful?

\bigskip

Would it be clearer to use two loops?

\bigskip

How about a version with single static assignments and slices? Without
slices?

\end{frame}

% ------------------------------------------------------------------------
%
\begin{frame}{Naïve search}

The input is a pattern~\(p\) of length~\(m\) and a text~\(t\) of
length~\(n\)

\bigskip

What configuration of the input leads to the minimal number of
comparisons? What is that number?

\bigskip

In the worst case, \(p\) is not in \(t\) and the test at line \(3\)
always fails on the last letter of \(p\), leading to make \(p\)
slide of one position at line \(5\) the latest as possible.

\bigskip

An example of positive match in the worst case is \(p =
\texttt{a}^{m-1}\texttt{b}\) and \(t = \texttt{a}^{n-1}\texttt{b}\),
with \(m \leq n\).

\bigskip

There are \(n - m + 1\) positions in \(t\) to compare with the first
letter of \(p\), and since there are \(m\) letters in \(p\), it makes
a total of \((n - m + 1) m = nm - m^2 + m < \len{t} \times \len{p}\)
positions.

\bigskip

What about the worst case when there is no match?

\end{frame}
